\documentclass[12pt,english]{article}

\usepackage[utf8]{inputenc}
\usepackage{amsmath}
\usepackage{cite}
\usepackage{graphicx}
\usepackage[a4paper,bindingoffset=0.2in,left=1.25in,right=.75in,top=1in, bottom = 1in, footskip=.25in]{geometry}
\usepackage{fourier}
\usepackage{inconsolata}
\usepackage{minted}
\usepackage{nameref}

\renewcommand{\baselinestretch}{1.25}
\graphicspath{ {../images/} }
\usemintedstyle{tango}

\begin{document}
\begin{center}
\includegraphics[scale = 0.5]{LincolnLogo}
\newline\newline
\thispagestyle{empty}
    \textbf{\huge{Classification of Biological Organisms from Images Using Advanced Mathematical Techniques} \\
    [100pt] \Large{Timothy Smith\\25796944} \\
    [40pt]\LARGE{Supervised by Dr Danilo Roccatano}\\
    [50pt]\today \\
	[40pt]Scientific Report in the 3\textsuperscript{rd} Year Module\\ MTH3009 Mathematics Project  }
\end{center}

\newpage
\section*{Abstract}

\newpage
\tableofcontents

\newpage
\section{Introduction}
\section{The Elliptic Fourier Transform}
The Fourier transform is a way of finding constituent frequencies within a function on the real domain.
The Fourier transform extends the concept of the Fourier series (which openerates on a bounded interval) to the real
domain.

\subsection{The Fourier Series}
The Fourier series
(\ref{eq:fourierseries})
lets you take a single repeating unit of a
periodic function and use this to generate an infinite sum of
sinusoidal functions that converge to it.
This is useful because trigonmetric functions have mathematical
properties that not all functions posess, allowing you to work
with them nicely in more situations.
\begin{equation}
	\label{eq:fourierseries}
	s(x) \sim A_0 + \sum_{n=1}^{\infty} \left(
		A_n \cos \left( \frac{2\pi nx}{P} \right) +
		B_n \sin \left( \frac{2\pi nx}{P} \right)
	\right)
\end{equation}
The \(\sim\) symbol indicates that the series doesn't strictly
converge in all cases.

Firstly we have a set of numbers, \(A_0\), \(A_n\) and \(B_n\):

\begin{equation}
\begin{aligned}
	A_0 &= \frac{1}{P} \int_{-\frac{P}{2}}^{\frac{P}{2}}
		s(x) \,dx \\
	A_n &= \frac{1}{P} \int_{-\frac{P}{2}}^{\frac{P}{2}}
		s(x) \cos \left(\frac{2\pi nx}{P}\right) \,dx \\
	A_b &= \frac{1}{P} \int_{-\frac{P}{2}}^{\frac{P}{2}}
		s(x) \sin \left(\frac{2\pi nx}{P}\right) \,dx \\
\end{aligned}
\end{equation}

Here \(A_0\) is a constant, and \(A_n\) and \(B_n\) are functions of \(n\).
In fact, since both the denominator of the fraction and the range of the integral for \(A_0\) are the same, \(P\),
\(A_0\) is simply the average around which the function oscillates.

A good intuition for why the Fourier series works is that the
integral of the product of a function and a sinusoid is a measure
of how well that function \textit{matches} the sinusoid.
For example, the $\sin$ function perfectly matches itself since 
$\sin^2x$ has a positive average, while the $\sin$ and $\cos$
functions are a perfect anti-match since $\sin x \cos x$ simplifies
to $\frac{\sin 2x}{2}$ which averages around $0$.

\subsection{The Fourier Transform}
The transform function ($S(t)$) for frequency $f$ is given by
(\ref{fourier:transform}).
\begin{equation} \label{fourier:transform}
	S(t)=\int_{-\infty}^{\infty}s(t)\cdot \mathrm{e}^{-i2\pi ft} dt
\end{equation}

\subsection{Fourier Analysis}

\subsection{Elliptical Fourier Analysis}
The seminal paper on this topic is \cite{kuhl:1982}.

For:
\begin{equation}
\begin{aligned}
	x(t) = A_0+\sum_{n=1}^{\infty}\left[a_n\cos\frac{2n\pi t}{T}+b_n\sin\frac{2n\pi t}{T}\right] \\
	y(t) = C_0+\sum_{n=1}^{\infty}\left[c_n\cos\frac{2n\pi t}{T}+d_n\sin\frac{2n\pi t}{T}\right]
\end{aligned}
\end{equation}

\be

This paper gives the coefficients (for a discrete set of points that form a closed contour):
\begin{equation}
\begin{aligned}
	a_n&=\frac{T}{2n^2\pi^2}\sum_{p=1}^{K}\left[\cos\frac{2n\pi t_p}{T}-\cos\frac{2n\pi t_{p-1}}{T}\right] \\
	b_n&=\frac{T}{2n^2\pi^2}\sum_{p=1}^{K}\left[\sin\frac{2n\pi t_p}{T}-\sin\frac{2n\pi t_{p-1}}{T}\right] \\
	c_n&=\frac{T}{2n^2\pi^2}\sum_{p=1}^{K}\left[\cos\frac{2n\pi t_p}{T}-\cos\frac{2n\pi t_{p-1}}{T}\right] \\
	d_n&=\frac{T}{2n^2\pi^2}\sum_{p=1}^{K}\left[\sin\frac{2n\pi t_p}{T}-\sin\frac{2n\pi t_{p-1}}{T}\right]
\end{aligned}
\end{equation}

It also gives the centres:

\begin{equation}
\begin{aligned}
	A_0 = \frac{1}{T}\sum_{p=1}^{K}\left[\frac{\Delta x_p}{2\Delta t_p}\left(t_p^2-t_{p-1}^2\right)+\xi_p\left(t_p-t_{p-1}\right) \right]\\
	C_0 = \frac{1}{T}\sum_{p=1}^{K}\left[\frac{\Delta y_p}{2\Delta t_p}\left(t_p^2-t_{p-1}^2\right)+\delta_p\left(t_p-t_{p-1}\right) \right]
\end{aligned}
\end{equation}

Where:

\begin{equation}
\begin{aligned}
	\xi_p = \sum_{j=1}^{p-1}\Delta x_j-\frac{\Delta x_p}{\Delta t_p}\sum_{j=1}^{p-1}\Delta t_j \\
	\delta_p = \sum_{j=1}^{p-1}\Delta y_j-\frac{\Delta y_p}{\Delta t_p}\sum_{j=1}^{p-1}\Delta t_j
\end{aligned}
\end{equation}

You can see this implemented in \nameref{code:grapharea}.

\section{Processing the Images}
In order to find shapes in our images for the purposes of analysis,
we will need to detect edges. For this we will use the Canny Edge
detection method, developed by John F. Canny.

\subsection{Convolution Kernels}
The most basic aspect of edge detection is the convolution kernel.
This is an $M{\times}M: \{M=2n+1, n\in\mathbb{N}\}$ matrix
(an odd sided square matrix).
Odd side lengths allow the kernel to be centered at each pixel.  If we let the function $f(x,y)$ represent the original image, $g(x,y)$ represent the convolved image and $\omega$
represent the kernel, then:
\begin{equation}
	g(x,y)=\sum_{i=-n}^{n}\sum_{j=-n}^{n}\omega(i,j)f(x-i,y-j)
\end{equation}
Notice that the coordinates of the kernel are not $1$ to $M$
as with a traditional matrix, but rather $-n$ to $n$.

The effect of this on an image will be to make each pixel of a
convolved image a function of the surrounding pixels.
The simplest convolution matrix is the identity convolution matrix,
(\ref{convm:identity}).
\begin{equation}
	\label{convm:identity}
	\begin{bmatrix}
		0 & 0 & 0\\
		0 & 1 & 0\\
		0 & 0 & 0\\
	\end{bmatrix}
\end{equation}
More examples include edge detection kernels like
(\ref{convm:edge1}) and (\ref{convm:edge2}).
\begin{center}
\begin{minipage}{.4\linewidth}
\begin{equation} \label{convm:edge1}
	\begin{bmatrix}
		 0 & -1 &  0\\
		-1 &  4 & -1\\
		 0 & -1 &  0\\
	\end{bmatrix}
\end{equation}
\end{minipage}
\begin{minipage}{.4\linewidth}
\begin{equation} \label{convm:edge2}
	\begin{bmatrix}
		-1 & -1 & -1\\
		-1 &  8 & -1\\
		-1 & -1 & -1\\
	\end{bmatrix}
\end{equation}
\end{minipage}
\end{center}

\subsection{Canny Edge Detection}

\section{Conclusion}

\section{Acknowledgements}

\section{Appendix}
Note: These are unlikely to be in the final report, they are simply there as an example appendix for now.
\subsection*{graph\_area.cpp}
\label{code:grapharea}
\inputminted[breakanywhere=true, tabsize=4]{cpp}{../code/elliptic_fourier/graph_area.cpp}

\section{Bibliography}
\bibliographystyle{ieeetr}
\bibliography{refs}

\end{document}
